\section{Implementation}\label{sec:implementation}

This chapter describes the implementation of the proposed system design.

\subsection{Technology stack}

The system is implemented using the following technologies:
\begin{itemize}
  \item Python 3.13 with NumPy~\cite{harris2020array} for the core processing logic;
  \item PostgreSQL 17 for data storage;
  \item Docker for containerized deployment;
  \item GitHub Actions for continuous integration.
\end{itemize}

\subsection{Architecture overview}

The implemented system follows the hybrid architecture selected in \cref{sec:design}.
The core module handles data processing and analysis, while extension modules provide specialized functionality.

% EXAMPLE: TikZ deployment diagram with several node shapes.
\begin{figure}[htbp]
  \centering
  \begin{tikzpicture}[
      alt={Deployment diagram: a client connects to a load balancer, which
           distributes requests to App Server 1 and App Server 2. Both
           servers use a shared cache in front of the database.},
      box/.style={rectangle, draw, minimum width=2.5cm, minimum height=0.8cm, align=center},
      client/.style={ellipse, draw, minimum width=2cm, minimum height=0.7cm, align=center},
      >=Stealth,
    ]
    \node[client] (client) {Client};
    \node[box, right=1cm of client] (lb) {Load Balancer};
    \node[box, above right=0.5cm and 1cm of lb] (app1) {App Server 1};
    \node[box, below right=0.5cm and 1cm of lb] (app2) {App Server 2};
    \node[box, right=4.5cm of lb] (cache) {Cache};
    \node[box, below=1.8cm of cache] (db) {Database};
    \draw[->] (client) -- (lb);
    \draw[->] (lb) -- (app1);
    \draw[->] (lb) -- (app2);
    \draw[->] (app1) -- (cache);
    \draw[->] (app2) -- (cache);
    \draw[->] (cache) -- (db);
  \end{tikzpicture}
  \caption{Implemented system architecture}\label{fig:impl-arch}
\end{figure}

% EXAMPLE: an external image from figures/. The alt text says what the
% image shows; the caption says why it is there.
\Cref{fig:sample-photo} demonstrates including an external image from the \path{figures/} directory.

\begin{figure}[htbp]
  \centering
  \includegraphics[width=0.4\linewidth, alt={Placeholder photograph used
    to demonstrate image inclusion}]{sample-photo}
  \caption{Sample image included from a file}\label{fig:sample-photo}
\end{figure}

% EXAMPLE: two images side by side in one figure, each with its own alt text.
\Cref{fig:two-images} shows two images placed side by side.
This is useful for comparing visual results or showing before and after states.

\begin{figure}[htbp]
  \centering
  \includegraphics[width=0.4\linewidth, alt={Placeholder photograph, state
    before processing}]{sample-photo}%
  \hfill
  \includegraphics[width=0.4\linewidth, alt={Placeholder photograph, state
    after processing}]{sample-photo}
  \caption{Two images side by side: before (left) and after (right)}\label{fig:two-images}
\end{figure}

\subsection{Key components}

The main processing module is responsible for data ingestion and transformation.
\Cref{lst:processor} shows a simplified version of the data processor.

% EXAMPLE: code listing. A "codelisting" float gives it a number and a
% caption; the code itself goes into a verbatim environment.
\begin{codelisting}[htbp]
  \caption{Data processor module (processor.py)}\label{lst:processor}
\begin{verbatim}
import json
import logging
from dataclasses import dataclass

logger = logging.getLogger(__name__)


@dataclass
class DataPoint:
    """A single data measurement."""
    timestamp: str
    value: float
    label: str | None = None


class DataProcessor:
    """Processes and transforms input data."""

    def __init__(self, config: dict):
        self.scale = config.get("scale", 1.0)
        self.results: list[dict] = []

    def load(self, path: str) -> list[DataPoint]:
        with open(path) as f:
            return [DataPoint(**item) for item in json.load(f)]

    def process(self, data: list[DataPoint]) -> list[dict]:
        for point in data:
            self.results.append({
                "timestamp": point.timestamp,
                "value": point.value * self.scale,
                "label": point.label or "unlabeled",
            })
        logger.info("Processed %d points", len(self.results))
        return self.results
\end{verbatim}
\end{codelisting}

\subsection{Configuration management}

The system uses a configuration file to manage runtime parameters.
\Cref{lst:config} shows an example configuration.

\begin{codelisting}[htbp]
  \caption{Example configuration file (config.json)}\label{lst:config}
\begin{verbatim}
{
  "scale": 1.5,
  "output_dir": "./results",
  "log_level": "INFO",
  "database": { "host": "localhost", "port": 5432 },
  "processing": { "batch_size": 100, "max_retries": 3 }
}
\end{verbatim}
\end{codelisting}

\subsection{Database schema}

\Cref{lst:sql} shows the database schema used for storing experimental results.

\begin{codelisting}[htbp]
  \caption{Database schema (schema.sql)}\label{lst:sql}
\begin{verbatim}
CREATE TABLE experiments (
    id          SERIAL PRIMARY KEY,
    name        VARCHAR(255) NOT NULL,
    created_at  TIMESTAMP DEFAULT CURRENT_TIMESTAMP,
    status      VARCHAR(50) DEFAULT 'pending'
);

CREATE TABLE results (
    id            SERIAL PRIMARY KEY,
    experiment_id INTEGER REFERENCES experiments(id),
    metric_name   VARCHAR(100) NOT NULL,
    metric_value  DOUBLE PRECISION NOT NULL
);
\end{verbatim}
\end{codelisting}

\subsection{Containerized deployment}

\Cref{lst:dockerfile} shows the Dockerfile used for deployment; container orchestration follows the practices described in~\cite{burns2016borg}.

\begin{codelisting}[htbp]
  \caption{Application Dockerfile}\label{lst:dockerfile}
\begin{verbatim}
FROM python:3.13-slim
WORKDIR /app
COPY requirements.txt .
RUN pip install --no-cache-dir -r requirements.txt
COPY src/ src/
EXPOSE 8080
CMD ["python", "-m", "src.main"]
\end{verbatim}
\end{codelisting}

\subsection{Security architecture}

The system implements the security controls recommended in~\cite{nist2020sp80053}, as shown in \cref{fig:security}.

\begin{figure}[htbp]
  \centering
  \begin{tikzpicture}[
      alt={Security pipeline: requests pass a firewall, then authentication,
           then encryption; authentication also writes to an audit log.},
      box/.style={rectangle, draw, minimum width=2.5cm, minimum height=0.8cm, align=center},
      >=Stealth,
    ]
    \node[box, fill=blue!10] (fw) {Firewall};
    \node[box, right=1cm of fw, fill=green!10] (auth) {Authentication};
    \node[box, right=1cm of auth, fill=orange!10] (enc) {Encryption};
    \node[box, below=0.8cm of auth, fill=yellow!10] (audit) {Audit Logging};
    \draw[->] (fw) -- (auth);
    \draw[->] (auth) -- (enc);
    \draw[->] (auth) -- (audit);
  \end{tikzpicture}
  \caption{Security architecture overview}\label{fig:security}
\end{figure}
